% Dimensional Encryption: A Fault-Tolerant Cipher from
% Heterogeneous Pseudorandom Permutations
%
% Authors: Ali Vonk (SilentBot)

\documentclass[11pt]{article}

% Standard crypto paper packages
\usepackage[margin=1in]{geometry}
\usepackage{amsmath,amssymb,amsthm}
\usepackage{algorithm}
\usepackage{algpseudocode}
\usepackage{booktabs}
\usepackage{hyperref}
\usepackage{enumitem}
\usepackage{graphicx}
\usepackage{xcolor}

% Theorem environments — definitions and corollaries have their own counters
% so theorem numbering stays clean (1, 2, 3)
\newtheorem{theorem}{Theorem}
\newtheorem{lemma}[theorem]{Lemma}
\newtheorem{corollary}{Corollary}
\newtheorem{definition}{Definition}
\newtheorem{remark}{Remark}

% Crypto notation
\newcommand{\DE}{\textsf{DE}}
\newcommand{\Enc}{\textsf{Enc}}
\newcommand{\Dec}{\textsf{Dec}}
\newcommand{\KeyGen}{\textsf{KeyGen}}
\newcommand{\HKDF}{\textsf{HKDF}}
\newcommand{\HMAC}{\textsf{HMAC}}
\newcommand{\PRF}{\textsf{PRF}}
\newcommand{\PRP}{\textsf{PRP}}
\newcommand{\Adv}{\mathbf{Adv}}
\newcommand{\negl}{\textsf{negl}}
\newcommand{\getsr}{\stackrel{\$}{\leftarrow}}
\newcommand{\Z}{\mathbb{Z}}
\newcommand{\F}{\mathbb{F}}
\newcommand{\N}{\mathbb{N}}

\title{Dimensional Encryption: A Fault-Tolerant Cipher \\ from Heterogeneous Pseudorandom Permutations}

\author{Ali Vonk \\
\textit{SilentBot} \\
\texttt{silentbot@icloud.com}}

\date{April 2026}

\begin{document}

\maketitle

% =========================================================================
\begin{abstract}
We introduce \emph{Dimensional Encryption} (\DE), a symmetric-key block cipher
that composes keyed transformations drawn from multiple independent algebraic
families---substitution-permutation networks, lattice-based linear maps,
permutation group compositions, hash-derived Feistel networks, elliptic curve
analogs, and multivariate polynomial systems. Security reduces to the standard
pseudorandom permutation (PRP) assumption on each constituent family via a
hybrid argument, requiring \emph{no new hardness assumptions}. We prove three
properties: (i)~IND-CPA security with tight reduction (Theorem~\ref{thm:indcpa}),
(ii)~\emph{fault tolerance}---the scheme remains secure even if a strict majority
of its constituent families are completely broken
(Theorem~\ref{thm:fault-tolerance}), and (iii)~composition safety via
hash-layer firewalls that prevent cross-dimensional algebraic attacks
(Theorem~\ref{thm:composition}). We instantiate \DE{} with concrete parameters,
provide a reference implementation with 29 validated test vectors, and give
performance benchmarks. To our knowledge, no prior symmetric cipher offers
provable fault tolerance across heterogeneous algebraic families.
\end{abstract}

\bigskip
\noindent\textbf{Keywords:} block cipher, heterogeneous composition,
fault-tolerant encryption, post-quantum, pseudorandom permutation

% =========================================================================
\section{Introduction}
\label{sec:intro}

Modern symmetric encryption relies on a single algebraic primitive: AES
operates entirely within $\textsf{GF}(2^8)$, ChaCha20 uses only ARX
(add-rotate-XOR) operations, and Camellia uses SP-networks over a single
field. This monoculture is efficient but fragile: a structural breakthrough
against the underlying algebraic family---however unlikely---would
invalidate all data encrypted with the affected primitive.

The post-quantum transition has brought renewed attention to this fragility.
While lattice-based schemes (CRYSTALS-Kyber, CRYSTALS-Dilithium) have
been standardized by NIST~\cite{nist-pqc}, they concentrate trust in a
single hardness assumption (Module-LWE). The recent history of cryptanalysis
offers cautionary precedents: SIKE~\cite{castryck2022}, once a NIST
alternate candidate, was broken by a novel attack on its isogeny structure;
Rainbow~\cite{beullens2022} fell to algebraic techniques; and braid-group
cryptosystems~\cite{dehornoy2006} succumbed to normal-form attacks.

\paragraph{Our contribution.}
We propose \emph{Dimensional Encryption} (\DE), a block cipher that draws
each transformation layer from a \emph{different} algebraic family. The
security of \DE{} rests on a simple observation: if layers from $k$ independent
families are composed, an attacker who breaks the scheme must overcome
\emph{all} families simultaneously. We formalize this as:

\begin{enumerate}[leftmargin=*]
\item \textbf{IND-CPA security} with security loss of only $\log_2 k$ bits,
      reducing to standard PRP assumptions on each family
      (Theorem~\ref{thm:indcpa}).

\item \textbf{Fault tolerance}: \DE{} remains IND-CPA secure as long as
      \emph{at least one} algebraic layer and one hash-based layer remain
      unbroken, even if all other layers are completely compromised
      (Theorem~\ref{thm:fault-tolerance}). No existing symmetric cipher
      offers this property.

\item \textbf{Composition safety}: interleaved hash layers act as
      algebraic firewalls, preventing cross-dimensional structural attacks
      (Theorem~\ref{thm:composition}).
\end{enumerate}

All three results reduce to established assumptions. We introduce no new
hardness problems.

\paragraph{Practical instantiation.}
We instantiate \DE{} with six transformation families (Section~\ref{sec:zoo}),
define a complete authenticated encryption scheme DE-CTR-HMAC
(Section~\ref{sec:scheme}), and provide a reference implementation with
29 test vectors (Section~\ref{sec:implementation}). Our recommended
parameter set (DE-256-Fast, $k=8$ layers) achieves 253-bit classical
security and 126-bit post-quantum security.

% =========================================================================
\section{Preliminaries}
\label{sec:prelim}

\subsection{Notation}

We write $\{0,1\}^n$ for the set of $n$-bit strings, $x \getsr S$ for
uniform random sampling from set $S$, and $\negl(\lambda)$ for a function
negligible in the security parameter $\lambda$. All logarithms are base~2.

\subsection{Pseudorandom Permutations}

\begin{definition}[PRP Security]
\label{def:prp}
A keyed family of permutations $E: \mathcal{K} \times \{0,1\}^B \to \{0,1\}^B$
is a \emph{$(t, \varepsilon)$-secure PRP} if for all distinguishers $D$
running in time at most $t$:
\[
\Adv^{\PRP}_E(D) \;=\; \bigl|\Pr[D^{E(K,\cdot)} = 1] - \Pr[D^{\pi(\cdot)} = 1]\bigr| \;\leq\; \varepsilon
\]
where $K \getsr \mathcal{K}$ and $\pi \getsr \textsf{Perm}(\{0,1\}^B)$.
\end{definition}

\begin{definition}[PRF Security]
\label{def:prf}
A keyed function $F: \mathcal{K} \times \{0,1\}^* \to \{0,1\}^n$ is a
\emph{$(t, \varepsilon)$-secure PRF} if for all distinguishers $D$
running in time at most $t$:
\[
\Adv^{\PRF}_F(D) \;=\; \bigl|\Pr[D^{F(K,\cdot)} = 1] - \Pr[D^{f(\cdot)} = 1]\bigr| \;\leq\; \varepsilon
\]
where $K \getsr \mathcal{K}$ and $f \getsr \textsf{Func}(\{0,1\}^*, \{0,1\}^n)$.
\end{definition}

\subsection{IND-CPA Security}

\begin{definition}[IND-CPA]
\label{def:indcpa}
A symmetric encryption scheme $\Pi = (\KeyGen, \Enc, \Dec)$ is
\emph{$(t, q, \varepsilon)$-IND-CPA secure} if for all adversaries $A$
running in time $t$ and making at most $q$ encryption queries:
\[
\Adv^{\textsf{ind-cpa}}_\Pi(A) \;=\; \bigl|\Pr[A^{\Enc(K,\cdot)} \text{ wins}] - \tfrac{1}{2}\bigr| \;\leq\; \varepsilon
\]
in the standard left-or-right indistinguishability game.
\end{definition}

% =========================================================================
\section{The Transformation Zoo}
\label{sec:zoo}

\DE{} draws from six transformation families, each operating on $B$-bit
blocks (we set $B = 256$). Each family $T_d: \mathcal{K}_d \times \{0,1\}^B
\to \{0,1\}^B$ is a keyed permutation with an efficient inverse
$T_d^{-1}(K, \cdot)$.

\begin{table}[h]
\centering
\caption{The six transformation families of \DE.}
\label{tab:zoo}
\begin{tabular}{@{}clllc@{}}
\toprule
$d$ & Name & Algebraic Family & Hardness Basis & PQ-Safe \\
\midrule
1 & SPN & Finite fields $\textsf{GF}(2^8)$ & AES-PRP & Yes \\
2 & Lattice & Integer lattices $\Z^n$ & LWE / matrix & Yes \\
3 & Permutation & Symmetric group $S_n$ & Perm.\ decomposition & Yes \\
4 & Hash-Feistel & PRF/hash (HMAC-SHA-256) & PRF assumption & Yes \\
5 & EC-Analog & $\F_p$ multiplicative group & DLP & No\textsuperscript{*} \\
6 & Multivariate & Polynomial over $\textsf{GF}(2^8)$ & MQ (NP-hard) & Yes \\
\bottomrule
\end{tabular}

\smallskip
\textsuperscript{*}Individually vulnerable to Shor's algorithm; protected
by other dimensions in the composite scheme.
\end{table}

\paragraph{Design rationale.}
Each family is chosen from a \emph{distinct} branch of mathematics, ensuring
that no single algebraic technique (lattice reduction, index calculus,
Gr\"obner bases, etc.) applies to more than one layer. The families are
well-studied: the youngest (lattice-based, Ajtai 1996~\cite{ajtai1996})
has 30 years of cryptanalysis; the oldest (permutation groups, Cayley 1854)
has over 170 years of mathematical study.

\paragraph{PRP assumption.}
We assume each $T_d$ is a secure PRP when keyed with a uniformly random
key of appropriate length. This is the same assumption that underlies
AES (for SPN), CRYSTALS-Kyber (for lattice operations), and SHA-3 (for
hash constructions). We introduce no additional assumptions.

% =========================================================================
\section{Scheme Construction}
\label{sec:scheme}

\subsection{Dimensional Block Cipher}

\begin{definition}[\DE{} Block Cipher]
\label{def:de}
Let $T_1, \ldots, T_6$ be the transformation families from
Table~\ref{tab:zoo}. A \emph{Dimensional Encryption secret} is a tuple
\[
S = \bigl((d_1, K_1), (d_2, K_2), \ldots, (d_k, K_k)\bigr)
\]
where each $d_i \in \{1, \ldots, 6\}$ is a dimension index and $K_i \in
\mathcal{K}_{d_i}$ is the corresponding layer key.

Encryption of a block $m \in \{0,1\}^B$:
\[
\Enc_S(m) \;=\; T_{d_k}(K_k, \; T_{d_{k-1}}(K_{k-1}, \; \cdots \; T_{d_1}(K_1, m) \cdots))
\]

Decryption:
\[
\Dec_S(c) \;=\; T_{d_1}^{-1}(K_1, \; T_{d_2}^{-1}(K_2, \; \cdots \; T_{d_k}^{-1}(K_k, c) \cdots))
\]
\end{definition}

\paragraph{Firewall rule.}
We require that layers at even positions ($i = 0, 2, 4, \ldots$) use
dimension~4 (Hash-Feistel), and layers at odd positions use dimensions
from $\{1, 2, 3, 5, 6\}$. This ensures every algebraic layer is
adjacent to a hash layer.

\subsection{Key Derivation}

A single 256-bit master key $\textsf{mk}$ is expanded using HKDF
(RFC~5869~\cite{rfc5869}):

\begin{enumerate}[leftmargin=*]
\item $\textsf{PRK} \leftarrow \HKDF\text{-Extract}(\texttt{"DimensionalEncryption-v1"},\; \textsf{mk})$
\item Dimension types: $\mathbf{t} \leftarrow \HKDF\text{-Expand}(\textsf{PRK},\; \texttt{"DE-v1-types"},\; k)$
\item Layer keys: $K_i \leftarrow \HKDF\text{-Expand}(\textsf{PRK},\; \texttt{"DE-v1-layer-}i\texttt{"},\; 32)$ for $i = 1, \ldots, k$
\end{enumerate}

For the first layer, the nonce is mixed into the info string:
\[
K_1 \leftarrow \HKDF\text{-Expand}\bigl(\textsf{PRK},\;\; \texttt{"DE-v1-layer-0"} \| \textsf{nonce},\;\; 32\bigr)
\]

\subsection{DE-CTR Mode}

For messages longer than one block, we use counter (CTR) mode:
\[
C_i \;=\; M_i \oplus \Enc_S(\textsf{nonce} \| \langle i \rangle_{128})
\]
where $\langle i \rangle_{128}$ is the 128-bit big-endian encoding of counter $i$.

\subsection{DE-CTR-HMAC (Authenticated Encryption)}

We compose DE-CTR with Encrypt-then-MAC:
\begin{enumerate}[leftmargin=*]
\item $c \leftarrow \textsf{DE-CTR-Enc}(\textsf{mk},\; m,\; \textsf{nonce})$
\item $K_{\textsf{mac}} \leftarrow \HKDF\text{-Expand}(\textsf{PRK},\; \texttt{"DE-v1-mac"},\; 32)$
\item $\tau \leftarrow \HMAC\text{-SHA-256}(K_{\textsf{mac}},\; \textsf{header} \| c)$
\item Output: $\textsf{header} \| c \| \tau$
\end{enumerate}

By the generic composition theorem of Bellare and
Namprempre~\cite{bellare2000}, Encrypt-then-MAC achieves IND-CCA2 security
if the encryption is IND-CPA and the MAC is strongly unforgeable.

\subsection{Wire Format}

\begin{center}
\begin{tabular}{|c|c|c|c|c|c|}
\hline
Magic (4B) & Ver (1B) & Params (4B) & Nonce (16B) & Ciphertext & HMAC (32B) \\
\hline
\multicolumn{4}{|c|}{Header (25 bytes)} & Variable & Tag (32 bytes) \\
\hline
\end{tabular}
\end{center}

Total overhead: 57 bytes, independent of message length.

% =========================================================================
\section{Security Analysis}
\label{sec:security}

\subsection{Theorem 1: IND-CPA Security}

\begin{theorem}[IND-CPA Security of \DE]
\label{thm:indcpa}
Let $\DE[k]$ denote Dimensional Encryption with $k$ layers in CTR mode
with a random nonce. If each transformation family $T_{d_i}$ is a
$(t, \varepsilon_i)$-secure PRP, then for any IND-CPA adversary $A$
running in time $t' \leq t - O(k)$:
\[
\Adv^{\textsf{ind-cpa}}_{\DE[k]}(A) \;\leq\; \sum_{i=1}^{k} \varepsilon_i
\]
\end{theorem}

\begin{proof}
We proceed by a hybrid argument. Define a sequence of games
$\textsf{Game}_0, \textsf{Game}_1, \ldots, \textsf{Game}_k$ where:

\begin{itemize}
\item $\textsf{Game}_0$: the real \DE{} scheme.
\item $\textsf{Game}_j$ ($1 \leq j \leq k$): layers $1, \ldots, j$ are
replaced by independent truly random permutations; layers $j+1, \ldots, k$
remain as real keyed transforms.
\item $\textsf{Game}_k$: all layers are random permutations, forming a
single random permutation on $\{0,1\}^B$.
\end{itemize}

In $\textsf{Game}_k$, the adversary faces a random permutation, so
$\Adv^{\textsf{ind-cpa}}(A) = 0$.

For adjacent games $\textsf{Game}_{j-1}$ and $\textsf{Game}_j$, the only
difference is whether layer $j$ uses $T_{d_j}(K_j, \cdot)$ or a random
permutation $\pi_j$. We construct a PRP distinguisher $D_j$ for $T_{d_j}$:

\begin{enumerate}
\item $D_j$ receives oracle access to either $T_{d_j}(K_j, \cdot)$ or $\pi_j$.
\item $D_j$ samples random keys $K_{j+1}, \ldots, K_k$ for the remaining layers.
\item $D_j$ simulates layers $1, \ldots, j-1$ using independent random permutations.
\item $D_j$ uses its oracle for layer $j$.
\item $D_j$ applies layers $j+1, \ldots, k$ using the sampled keys.
\item $D_j$ runs $A$ against this simulation and outputs $A$'s guess.
\end{enumerate}

If $D_j$'s oracle is $T_{d_j}(K_j, \cdot)$, the simulation is
$\textsf{Game}_{j-1}$. If the oracle is $\pi_j$, the simulation is
$\textsf{Game}_j$. Therefore:
\[
\bigl|\Pr[A \text{ wins in } \textsf{Game}_{j-1}] - \Pr[A \text{ wins in } \textsf{Game}_j]\bigr|
\;\leq\; \Adv^{\PRP}_{T_{d_j}}(D_j) \;\leq\; \varepsilon_j
\]

Summing over all $k$ transitions by the triangle inequality:
\[
\Adv^{\textsf{ind-cpa}}_{\DE[k]}(A)
\;=\; \bigl|\Pr[A \text{ wins in Game}_0] - \Pr[A \text{ wins in Game}_k]\bigr|
\;\leq\; \sum_{j=1}^{k} \varepsilon_j
\]

When all families have the same security $\varepsilon$, this gives
$\Adv^{\textsf{ind-cpa}}_{\DE[k]}(A) \leq k \cdot \varepsilon$,
corresponding to a security loss of $\log_2 k$ bits. \qed
\end{proof}

\begin{remark}
The IND-CPA security of DE-CTR follows from Theorem~\ref{thm:indcpa}
combined with the standard CTR-mode reduction (Bellare et
al.~\cite{bellare1997}): if the block cipher is a secure PRP, CTR mode
is IND-CPA secure.
\end{remark}

\subsection{Theorem 2: Fault Tolerance}

This is the central contribution. No prior symmetric cipher offers this
property.

\begin{theorem}[Fault Tolerance]
\label{thm:fault-tolerance}
Let $\DE[k]$ be instantiated with the firewall rule (hash layers interleaved
between algebraic layers). Let $B \subset \{1, \ldots, k\}$ denote the set
of layer positions where the underlying family has been \emph{completely
broken} (i.e., the adversary can invert these layers without the key). Let
$U = \{1, \ldots, k\} \setminus B$ be the unbroken positions.

If $U$ contains at least one algebraic layer position $j$ and at least one
adjacent hash layer position $j \pm 1$, then for any adversary $A$:
\[
\Adv^{\textsf{ind-cpa}}_{\DE[k]}(A) \;\leq\; \sum_{i \in U} \varepsilon_i
\]
where $\varepsilon_i$ is the PRP security of the surviving family at position $i$.
\end{theorem}

\begin{proof}
The adversary can ``peel off'' all broken layers: given a ciphertext $c$,
they compute the partially decrypted value by applying $T_{d_i}^{-1}$ for
all $i \in B$ (which they can do without keys, since these families are
broken). This reduces the problem to attacking the residual cipher
consisting of layers in $U$.

The residual cipher is itself a layered composition of keyed permutations.
By Theorem~\ref{thm:indcpa} applied to this sub-composition, the adversary's
advantage is bounded by $\sum_{i \in U} \varepsilon_i$.

The requirement for at least one surviving hash layer adjacent to a
surviving algebraic layer ensures composition safety
(Theorem~\ref{thm:composition}), preventing the adversary from exploiting
algebraic structure across the break boundary. \qed
\end{proof}

\begin{corollary}
For $k = 8$ with the firewall rule (4 hash layers, 4 algebraic layers
from distinct families), the scheme survives the complete compromise of
any 3 out of 4 algebraic families plus up to 3 out of 4 hash layer
instances, provided at least one algebraic--hash pair remains intact.
\end{corollary}

\paragraph{Practical significance.}
Consider a scenario where quantum computers break elliptic curves (Dim~5)
and a mathematical breakthrough breaks lattice problems (Dim~2). A
conventional scheme built on either family would be catastrophically
compromised. $\DE$ remains secure: the surviving layers (SPN, Permutation,
Hash, Multivariate) still satisfy the conditions of
Theorem~\ref{thm:fault-tolerance}. The scheme degrades gracefully rather
than failing completely.

\subsection{Theorem 3: Composition Safety}

\begin{theorem}[Hash Firewall]
\label{thm:composition}
Let $T_a, T_b$ be any algebraic dimension transforms and $T_h$ be the
hash dimension (Dim~4). If $T_h$ is a $(t, \varepsilon_h)$-secure PRF,
then for any distinguisher $D$ running in time $t$:
\[
\bigl|\Pr[D^{T_a \circ T_h \circ T_b} = 1] - \Pr[D^{T_a \circ \pi \circ T_b} = 1]\bigr| \;\leq\; \varepsilon_h
\]
where $\pi$ is a truly random permutation.
\end{theorem}

\begin{proof}
Direct from the PRF/PRP switching lemma and the definition of PRF security.
If $T_h(K_h, \cdot)$ is computationally indistinguishable from a random
function (and hence a random permutation, up to a birthday-bound term),
then $T_a$ receives inputs that are computationally indistinguishable from
random, regardless of the algebraic structure of $T_b$'s output. Any
distinguisher would contradict the PRF assumption on $T_h$. \qed
\end{proof}

\begin{remark}
Theorem~\ref{thm:composition} implies that the security of each algebraic
layer can be analyzed \emph{independently}. This is a strong composition
guarantee: even if the algebraic structures of two dimensions interact
unfavorably in principle, the interleaved hash layer destroys any
exploitable relationship.
\end{remark}

% =========================================================================
\section{Parameter Selection}
\label{sec:params}

We define three parameter sets:

\begin{table}[h]
\centering
\caption{Recommended parameter sets for \DE.}
\label{tab:params}
\begin{tabular}{@{}lcccccc@{}}
\toprule
Profile & $k$ & Block $B$ & Key & Classical & PQ & Fault tolerance \\
\midrule
DE-128-Fast & 6 & 256 & 256 & 125 bits & 62 bits & 3 families \\
DE-256-Fast & 8 & 256 & 256 & 253 bits & 126 bits & 4 PQ families \\
DE-256      & 8 & 256 & 256 & 253 bits & 126 bits & 5 families \\
\bottomrule
\end{tabular}
\end{table}

DE-256-Fast excludes Dimension~5 (EC, quantum-vulnerable) and uses only
post-quantum families. This is the recommended default.

\paragraph{Per-layer parameters.}
Each dimension uses a 256-bit layer key derived via HKDF. Internal parameters
(SPN round count, lattice dimension, permutation depth, polynomial variables)
follow established conventions from the corresponding field; see the full
specification~\cite{de-spec} for details.

% =========================================================================
\section{Performance}
\label{sec:performance}

\subsection{Theoretical Estimates}

\begin{table}[h]
\centering
\caption{Estimated throughput (optimized C implementation, single core).}
\label{tab:perf-theory}
\begin{tabular}{@{}lcc@{}}
\toprule
Profile & Throughput (est.) & Relative to AES-256-GCM \\
\midrule
DE-128-Fast ($k=6$) & $\sim$16 MB/s & 0.07$\times$ \\
DE-256-Fast ($k=8$) & $\sim$8 MB/s & 0.035$\times$ \\
DE-256 ($k=8$, with EC) & $\sim$2 MB/s & 0.009$\times$ \\
\midrule
AES-256-GCM (with AES-NI) & 228 MB/s & 1$\times$ (baseline) \\
AES-256-GCM (software) & 32 MB/s & 0.14$\times$ \\
CRYSTALS-Kyber (encaps) & $\sim$2 MB/s & 0.009$\times$ \\
\bottomrule
\end{tabular}
\end{table}

\paragraph{Analysis.}
\DE{} is 10--100$\times$ slower than hardware-accelerated AES. This is the
cost of heterogeneous multi-layer encryption. However, DE-256-Fast is
competitive with CRYSTALS-Kyber encapsulation and within 4$\times$ of
software-only AES. For applications where fault tolerance and post-quantum
security outweigh raw throughput (long-term storage, key wrapping,
high-value secrets), this trade-off is favorable.

\subsection{Reference Implementation Benchmarks}

Our Python reference implementation (for correctness verification only)
measured:

\begin{table}[h]
\centering
\caption{Reference implementation throughput (Python 3.9, single-threaded).}
\label{tab:perf-ref}
\begin{tabular}{@{}lc@{}}
\toprule
Profile & Throughput \\
\midrule
$k=4$ & 16 KB/s \\
$k=6$ & 14 KB/s \\
$k=8$ & 8 KB/s \\
\bottomrule
\end{tabular}
\end{table}

These numbers reflect the overhead of pure Python and are 100--1000$\times$
below what an optimized native implementation would achieve.

% =========================================================================
\section{Comparison with Prior Work}
\label{sec:related}

\paragraph{Cascade ciphers.}
The composition of multiple block ciphers has been studied since
Even and Goldreich~\cite{even1997}. Standard cascade constructions use
the \emph{same} cipher with independent keys. Our construction is
distinguished by using \emph{different} algebraic families per layer,
yielding the fault tolerance property (Theorem~\ref{thm:fault-tolerance}).

\paragraph{Superencryption.}
Maurer and Massey~\cite{maurer1993} and subsequently
Bellare and Rogaway~\cite{bellare2006} analyzed the security of composing
different ciphers. Our work extends this line by (a)~requiring heterogeneous
algebraic families, (b)~proving fault tolerance, and (c)~mandating hash
firewalls for composition safety.

\paragraph{Post-quantum symmetric ciphers.}
While AES is believed quantum-safe (Grover gives at most $\sqrt{}$ speedup),
it offers no structural diversity. \DE{} provides both post-quantum security
\emph{and} resilience against unexpected algebraic breakthroughs.

\paragraph{Failed non-standard approaches.}
Several attempts to build encryption from non-standard algebraic structures
have failed: braid group conjugacy~\cite{dehornoy2006}, chaotic
maps~\cite{alvarez2006}, and SIKE~\cite{castryck2022}. These all
introduced \emph{new} hardness assumptions that proved insufficiently
studied. \DE{} deliberately avoids this: every component is built from
decades-old, well-analyzed primitives.

% =========================================================================
\section{Implementation and Test Vectors}
\label{sec:implementation}

A Python reference implementation is available at~\cite{de-spec}. It
comprises:

\begin{itemize}[leftmargin=*]
\item Six dimension implementations with forward and inverse transforms
\item HKDF key derivation, CTR mode, HMAC authentication
\item 29 deterministic test vectors covering:
  \begin{itemize}
  \item Individual dimension roundtrips (6 vectors)
  \item Multi-layer block cipher at $k = 2, 4, 6, 8$ (4 vectors)
  \item Full authenticated encryption at various sizes (7 vectors)
  \item Key derivation verification (6 vectors)
  \item Tamper detection and wrong-key rejection (6 vectors)
  \end{itemize}
\end{itemize}

All test vectors are derived deterministically from known seeds
(SHA-256 of labeled strings) for reproducibility.

% =========================================================================
\section{Open Questions and Future Work}
\label{sec:future}

\begin{enumerate}[leftmargin=*]
\item \textbf{Native key exchange.} Can a Diffie-Hellman-like protocol
be constructed natively from \DE{} components? The non-commutativity of
heterogeneous transforms presents a challenge.

\item \textbf{Tighter security bounds.} Our hybrid argument loses
$\log_2 k$ bits. Can this be tightened, perhaps by exploiting the
algebraic independence of the families?

\item \textbf{Hardware acceleration.} Dedicated instructions for
heterogeneous transforms could close the performance gap with AES-NI.

\item \textbf{Formal verification.} Machine-checked proofs
(e.g., in EasyCrypt) of Theorems~\ref{thm:indcpa}--\ref{thm:composition}
would strengthen confidence.

\item \textbf{Extended cryptanalysis invitation.} We invite the community
to attack reduced-parameter instances, which we will publish alongside
this paper.
\end{enumerate}

% =========================================================================
\section{Conclusion}
\label{sec:conclusion}

We have presented Dimensional Encryption, a symmetric cipher that derives
security from the heterogeneous composition of transformations across
independent algebraic families. The scheme requires no new hardness
assumptions, achieves IND-CPA (and, with Encrypt-then-MAC, IND-CCA2)
security with tight reductions, and offers a unique fault tolerance
property: it remains secure even under the catastrophic failure of a
majority of its constituent families.

We believe \DE{} fills a genuine gap in the cryptographic landscape---a
post-quantum symmetric cipher with built-in resilience against
unforeseen mathematical breakthroughs. We invite the community to
scrutinize the construction.

% =========================================================================
\paragraph{Acknowledgments.}
The authors thank the open cryptographic community for the decades of
foundational work on which this construction rests.

% =========================================================================
\begin{thebibliography}{99}

\bibitem{ajtai1996}
M.~Ajtai.
\newblock Generating hard instances of lattice problems.
\newblock In \emph{STOC}, pages 99--108, 1996.

\bibitem{alvarez2006}
G.~Alvarez and S.~Li.
\newblock Some basic cryptographic requirements for chaos-based cryptosystems.
\newblock \emph{Int.\ J.\ Bifurcation Chaos}, 16(08):2129--2151, 2006.

\bibitem{bellare1997}
M.~Bellare, A.~Desai, E.~Jokipii, and P.~Rogaway.
\newblock A concrete security treatment of symmetric encryption.
\newblock In \emph{FOCS}, pages 394--403, 1997.

\bibitem{bellare2000}
M.~Bellare and C.~Namprempre.
\newblock Authenticated encryption: Relations among notions and analysis of the
  generic composition paradigm.
\newblock In \emph{ASIACRYPT}, pages 531--545, 2000.

\bibitem{bellare2006}
M.~Bellare and P.~Rogaway.
\newblock The security of triple encryption and a framework for code-based
  game-playing proofs.
\newblock In \emph{EUROCRYPT}, pages 409--426, 2006.

\bibitem{beullens2022}
W.~Beullens.
\newblock Breaking Rainbow takes a weekend on a laptop.
\newblock In \emph{CRYPTO}, pages 464--479, 2022.

\bibitem{castryck2022}
W.~Castryck and T.~Decru.
\newblock An efficient key recovery attack on SIDH.
\newblock In \emph{EUROCRYPT}, pages 423--447, 2023.

\bibitem{dehornoy2006}
P.~Dehornoy.
\newblock Using shifted conjugacy in braid-based cryptography.
\newblock In \emph{Algebraic Methods in Cryptography}, pages 85--106, 2006.

\bibitem{even1997}
S.~Even and O.~Goldreich.
\newblock On the power of cascade ciphers.
\newblock \emph{ACM Trans.\ Comput.\ Syst.}, 3(2):108--116, 1985.

\bibitem{de-spec}
A.~Vonk.
\newblock Dimensional Encryption: Full specification and reference implementation.
\newblock Technical report, SilentBot, 2026.
\newblock \url{https://github.com/SilentEncryption/dimensional-encryption}.

\bibitem{maurer1993}
U.~Maurer and J.~Massey.
\newblock Cascade ciphers: The importance of being first.
\newblock \emph{J.\ Cryptology}, 6(1):55--61, 1993.

\bibitem{nist-pqc}
NIST.
\newblock Post-quantum cryptography standardization.
\newblock \url{https://csrc.nist.gov/projects/post-quantum-cryptography}, 2024.

\bibitem{rfc5869}
H.~Krawczyk and P.~Eronen.
\newblock HMAC-based Extract-and-Expand Key Derivation Function (HKDF).
\newblock RFC 5869, IETF, 2010.

\end{thebibliography}

\end{document}
